This section presents plots for single instance based convergence and tables averaged over multiple instances.

In figure \ref{fig:alpha0.0} we can see the convergence to stationarity for the original matrix $P_0$ as defined on page \pageref{sec:P0_def}.
For the first two steps, the L1 distance decreases faster than the theoretical expectation of $\mu^t$ but from there on until the end of calculations when the distance is smaller than \num{1e-5} the slopes of the lines are identical, the the behaviour matches the theoretical expectation.


\begin{figure}[h]
	\centering
	\includegraphics[width=0.75\linewidth]{figs/plot_alpha0.0}
	\caption{Convergence Behaviour for Original Matrix $P_0$}
	\label{fig:alpha0.0}
\end{figure}
\newpage

In figure \ref{fig:alpha0.5}, we can see the convergence when the Markov chain is mixed with the identity matrix at half contribution from each.
Again, the actual convergence dips slightly below the asymptotic theoretical curve but runs within visible accuracy perfectly to the theoretical line from approximately $t>5$.
From the diagram it is also visible that it takes a bit more than twice as long to reach the final threshold.
\begin{figure}[h]
	\centering
	\includegraphics[width=0.75\linewidth]{figs/plot_alpha0.5}
	\caption{Convergence Behaviour for Lazy Matrix With $\alpha=0.5$}
	\label{fig:alpha0.5}
\end{figure}
\newpage

When further raising the laziness factor to $\alpha=0.8$, the convergence behaviour looks as in figure \ref{fig:alpha0.8}.
We can see the same dip and subsequent parallel behaviour as in the previous plots.
The convergence slowed down further and is then a factor of about 3 slower than for $\alpha=0.5$.

\begin{figure}[h]
	\centering
	\includegraphics[width=0.75\linewidth]{figs/plot_alpha0.8}
	\caption{Convergence Behaviour for Lazy Matrix With $\alpha=0.8$}
	\label{fig:alpha0.8}
\end{figure}
\newpage

In table \ref{tab:conv_x4=1} the results of a single convergence for multiple different lazy matrices are shown.
The starting condition was the same in all cases and the same as in the figures \ref{fig:alpha0.0},  \ref{fig:alpha0.5} and \ref{fig:alpha0.8}, so $x0 = (0, 0, 0, 1)^T$.
The table shows the inverse relation between $g$ and the convergence time quite well.
For the cases where $g$ is not particularly close to one it holds only approximately.
When going from $\alpha=0$ to $0.5$, $g$ is approximately halved while convergence slows down by a bit more than a factor of 2.
For $\alpha=0.8$, $g$ is reduced by a factor of five compared to the original matrix and the convergence time increases more than sixfold.
For $\alpha=0.9$, the spectral gap is halved and the convergence time doubled compared to the previous $\alpha=0.8$ case.
When adding one further nine, so going from $0.9$ to $0.99$ and from $0.99$ to $0.999$, the spectral gap is reduced by one order of magnitude and the convergence time increases by approximately the same factor.
Especially for the higher end of the laziness factors, the theoretically derived approximate proportionality between $\frac{1}{g}$ and the convergence time holds up under numerical experiments.

\begin{table}[h]
	\centering
\begin{tabular}{lccr}
	\hline
	Label & $\mu$ &$g$ &  Time to go below \num{0.001}\\
	\hline
	$P_{0.0}$ & 0.478078 &0.521922 &  9 \\
	$P_{0.5}$ & 0.739039 &0.260961 &  21 \\
	$P_{0.8}$ & 0.895616 &0.104384 &  58 \\
	$P_{0.9}$ & 0.947808 &0.052192 &  118 \\
	$P_{0.99}$ & 0.994781 &0.005219 &  1205 \\
	$P_{0.999}$ & 0.999478 &0.000522 &  12069 \\
\end{tabular}
\caption{Convergence Table for $x0 = (0, 0, 0, 1)^T$}
\label{tab:conv_x4=1}
\end{table}

Table \ref{tab:conv_stats} shows statistics of the convergence speed gathered using 10000 different starting distributions.
Each was randomly sampled by taking a vector of four independent random numbers in the unit interval and normalizing the result.
This lead to a mean initial variation distance of \num{0.23} which is between three and four times smaller than the variation distance of $(0, 0, 0, 1)^T$ which explains the smaller mean convergence time in table \ref{tab:conv_stats} compared to the previous table \ref{tab:conv_x4=1}.
Another issue is that the variation distance and the convergence time do not follow the same distribution which makes quantitative comparisons between the single instance and averaged case difficult.
However, this is not too relevant here because the main goal is to further support the theoretical finding that the convergence time is approximately inversely related to to spectral gap.
This finding is well replicated, similarly to the single instance case much more precisely at the upper end of the tested range for $\alpha$.
It is interesting to see how the standard deviation of the convergence behaves almost perfectly linear over the entire range of $\alpha$ values.  
\begin{table}[h]
	\centering
	\begin{tabular}{lccc}
		\hline
		Label &$g$ &  Mean Convergence Time $t_C$ & Standard Deviation of $t_C$ \\
		\hline
		$P_{0.0}$ &0.521922 &  6.6118 & 1.11548 \\
		$P_{0.5}$ &0.260961 & 15.7148 & 2.3642 \\
		$P_{0.8}$ &0.104384 & 42.5739 & 6.07746 \\
		$P_{0.9}$ &0.052192 & 87.3741 & 12.3557 \\
		$P_{0.99}$ &0.005219 & 892.331 & 124.477 \\
		$P_{0.999}$ &0.000522 & 8944.04 & 1260.39 \\
		\hline
	\end{tabular}
	\caption{Statistics for Convergence for the Different Matrices}
	\label{tab:conv_stats}
\end{table}